\chapter{Tuple 與 unpacking}
\label{ch:tuples}

前面幾章你已經學過怎麼用 list 把一堆東西裝在一起、怎麼用 dict 把「名字」對應到「值」。
這一章要認識另一種裝東西的容器，叫 tuple（唸作「thu-pl」或「too-pl」都有人講，沒有標準
答案）。

乍看之下 tuple 很像 list：一樣可以用中括號的方式取出裡面第幾個東西，一樣可以用
\texttt{for} 走訪，一樣可以串接，但它有一個關鍵的差別：\textbf{造出來之後，裡面的內容
就不能再改了}。這聽起來像是缺點（少了一個功能），但正是這個「不能改」的特性，讓 tuple
在很多情境下比 list 更好用，你會在這一章慢慢看到為什麼。
\index{tuple}

這一章分三段推進：

\begin{itemize}
\item 第一段先讓你認識 tuple 這個新容器本身：怎麼造一個 tuple、它跟 list
有哪些操作是共通的（包括串接與重複這兩個你在字串、list 章節都見過的運算子）、然後親手
撞一次「不能修改」這道牆，看看 Python 會丟出什麼樣的錯誤訊息，順便簡單認識一個以 tuple
為基礎、可以額外用名字取值的進階變體。
\item 第二段是這一章真正的重頭戲：把 tuple 跟賦值語法
結合在一起的技巧，統稱「unpacking」（解包）。一行程式碼同時把好幾個值分別存進好幾個變數
裡，這個技巧會讓你少寫很多先存成一個容器、再一個一個挑出來的繁瑣步驟，日常寫程式會一直
用到它，包括等一下你會學到的「函式同時回傳好幾個結果」也是靠它撐起來的，你還會學到用星
號一次接住「剩下的」那些值，以及用底線忽略用不到的值這兩個實用的小技巧。
\item 第三段把視野拉
高一層，正式回答一個從第 \ref{ch:variables} 章就埋下的問題：如果「等號是幫物件貼名字牌，不是複製」這件事
是真的，那為什麼 tuple 看起來「整個換掉」的時候，感覺就像被複製了一份新的一樣？這一段會
用具體的例子（外加一個能親眼看見「這是不是同一個物件」的小工具）證明，其實 tuple 背後遵
守的規則跟 list 一模一樣，只是因為它不能被原地修改，所以任何「看起來像修改」的操作，骨
子裡都只能是「造一個新物件，把名字牌重新貼上去」。
\end{itemize}

\section{認識 Tuple：先像 List，再撞見「不能改」}
\label{sec:tuple-basics}

要造一個 tuple，最常見的寫法是用一對圓括號把幾個值包起來，中間用逗號分隔；其實圓括號
不是必要的，逗號本身才是關鍵，但為了可讀性，大多數時候還是會加上括號。如果手上已經有一
個 list 或其他容器，也可以直接用 \texttt{tuple()} 把它轉換成 tuple。
\index{tuple()}

\begin{code}{python}{建立 tuple 的幾種寫法}
# 用括號建立（最常見的寫法）
t = (x, y, z)

# 逗號才是重點，括號其實可以省略
t = x, y, z

# 從其他序列轉換（例如 list）
t = tuple(some_container)
\end{code}

\begin{example}{python}{(make-coordinate.py) 建立一組座標}
latitude = 25.03
longitude = 121.56
coordinate = (latitude, longitude)
print(coordinate)
print(type(coordinate))

# 也可以把現有的 list 轉換成 tuple
scores_list = [88, 92, 76]
scores_tuple = tuple(scores_list)
print(scores_tuple)
\end{example}

\iosample{}{(25.03, 121.56) \\ <class 'tuple'> \\ (88, 92, 76) \\}

留意 \texttt{print(coordinate)} 印出來的樣子：Python 用圓括號顯示 tuple，這也是為什麼
大部分人習慣用括號來建立它，看起來跟印出來的樣子一致，比較好對照。

\subsection*{tuple 也支援你已經學過的操作}

你會發現，tuple 支援的操作你幾乎都學過，因為第 \ref{ch:lists} 章介紹 list 的時候，用索引取值、切
片、\texttt{in}、\texttt{len()}、\texttt{for} 走訪，這些操作在 tuple 上的寫法跟意義完
全一樣，不需要重新學一次。

\begin{example}{python}{(tuple-like-list.py) tuple 也能索引、切片、走訪}
colors = ("red", "green", "blue")

print(colors[0])          # 索引，跟 list 一樣
print(colors[1:])         # 切片，跟 list 一樣
print("green" in colors)  # 用 in 檢查成員關係
print(len(colors))        # 元素數量

for color in colors:      # for 迴圈走訪，跟 list 一樣
    print(color)
\end{example}

\iosample{}{red \\ ('green', 'blue') \\ True \\ 3 \\ red \\ green \\ blue \\}

\subsection*{tuple 也可以用 \texttt{+} 串接、用 \texttt{*} 重複}

除了上面這幾個操作，第 \ref{ch:strings} 章字串、第 \ref{ch:lists} 章 list 都教過另外兩個運算子，tuple 一樣支援：
用 \texttt{+} 把兩個 tuple 接成一個新的 tuple，用 \texttt{*} 把一個 tuple 的內容重複
好幾次。

\begin{code}{python}{tuple 的串接與重複}
combined = t1 + t2   # 把兩個 tuple 合併成一個新的 tuple
repeated = t * n      # 把 t 的內容重複 n 次
\end{code}

\begin{example}{python}{(tuple-concat-repeat.py) 串接時段、預留固定長度的佔位內容}
morning = ("09:00", "10:00", "11:00")
afternoon = ("14:00", "15:00")
schedule = morning + afternoon
print(schedule)

placeholder = ("-",) * 5
print(placeholder)
\end{example}

\iosample{}{('09:00', '10:00', '11:00', '14:00', '15:00') \\ ('-', '-', '-', '-', '-') \\}

這兩個運算子的行為跟第 \ref{ch:lists} 章 list 版本的 \texttt{+}、\texttt{*} 完全對應：\texttt{+}
不會修改任何一邊原本的 tuple，會造出一個全新的 tuple 把兩邊接起來；\texttt{*} 也不是
把每個元素乘上一個數字，是把整份內容重複貼上好幾次。\texttt{("-",) * 5} 這種寫法在需
要先準備好一份固定長度、內容之後才會確定的容器時很好用。

\subsection*{tuple 之間也可以直接比較大小}

tuple 除了可以用 \texttt{==} 比較兩個 tuple 的內容是否完全一樣，也可以直接用
\texttt{<}、\texttt{>} 這些比較運算子比較「大小」。比較的規則跟比對兩個字串很像：先比
第一個元素，如果分得出高下就直接得到答案，分不出來才往下比第二個、第三個，以此類推，這
種比法有個名字叫「字典順序比較」，因為概念上跟查字典時比較兩個詞先後順序的方式一樣。

\begin{example}{python}{(compare-tuples.py) 比較版本號、依完賽時間排序}
version_a = (1, 4, 2)
version_b = (1, 5, 0)
print(version_a < version_b)

runner_x = ("11:32", "Alice")
runner_y = ("11:28", "Bob")
print(sorted([runner_x, runner_y]))
\end{example}

\iosample{}{True \\{} [('11:28', 'Bob'), ('11:32', 'Alice')] \\}

第一個例子裡，\texttt{version\_a} 和 \texttt{version\_b} 的第一個元素都是
\texttt{1}，分不出高下，於是往下比第二個元素，\texttt{4 < 5}，比較到這裡就有答案了，
後面的元素完全不會再看。第二個例子把「完賽時間、姓名」包成 tuple 放進 list 裡，直接呼
叫第 \ref{ch:lists} 章學過的 \texttt{sorted()}，因為第一個元素（時間字串）就已經能分出順序，
\texttt{sorted()} 自然依完賽時間由先到後排好，這正是把好幾筆相關資料包成 tuple、再交
給排序函式處理的常見寫法。

\stylepoint{把好幾筆相關資料包成 tuple 放進 list 裡，再交給 \texttt{sorted()} 排序，是
很常見的資料處理寫法。因為 tuple 之間的比較會自動照順序一路比下去，你不用自己寫比較邏
輯，只要把最重要的排序依據放在 tuple 的第一個位置就好，這個小技巧之後處理表格型資料時
會很有用。}

\subsection*{撞牆：tuple 不能被修改}

既然 tuple 這麼像 list，你可能會想，那對它做 \texttt{append()} 或直接改某個位置的值，
應該也可以吧？試試看就知道了。

\begin{example}{python}{(change-tuple.py) 嘗試修改 tuple 的內容}
colors = ("red", "green", "blue")
colors[0] = "purple"
\end{example}

存檔執行後，Python 不會讓你得逞：

\begin{verbatim}
Traceback (most recent call last):
  File "change-tuple.py", line 2, in <module>
    colors[0] = "purple"
    ~~~~~~^^^
TypeError: 'tuple' object does not support item assignment
\end{verbatim}

\texttt{TypeError} 是說「這個型別不支援你要做的這個操作」；\texttt{does not support
item assignment} 說得更白：「不支援用索引賦值來改變裡面某一個位置的內容」。這句話你要
牢牢記住，這件事跟「變數到底是什麼」有關，我們留到這章最後一節再回頭講清楚。

那如果改用 \texttt{.append()} 呢？結果甚至不是同一種錯誤：tuple 根本沒有
\texttt{.append()} 這個方法，Python 連「有沒有這個操作」都判斷不出來，直接告訴你這個
方法不存在：

\begin{verbatim}
Traceback (most recent call last):
  File "change-tuple.py", line 2, in <module>
    colors.append("purple")
    ^^^^^^^^^^^^^
AttributeError: 'tuple' object has no attribute 'append'
\end{verbatim}

\keypoint{list 有 \texttt{.append()}、\texttt{.remove()}、\texttt{.sort()} 這些「會改
變自己」的方法，tuple 完全沒有這些方法。這不是「有但被鎖住」，是打從一開始就沒有設計這
些方法，因為 tuple 造出來之後本來就不打算被改變。}

\subsection*{單元素 tuple 的陷阱：別忘了逗號}

還有一個很容易踩到的坑：如果你只想放一個元素進 tuple，光加括號是不夠的。

\begin{example}{python}{(single-element-tuple.py) 忘記逗號的陷阱}
not_a_tuple = ("p")
print(type(not_a_tuple))

really_a_tuple = ("p",)
print(type(really_a_tuple))
\end{example}

\iosample{}{<class 'str'> \\ <class 'tuple'> \\}

原因是：圓括號在 Python 裡本來就有另一個用途，用來改變運算的優先順序（例如
\texttt{(1 + 2) * 3}），所以單獨一個 \texttt{("p")} 會被 Python 當成括號裡包著一個字
串，跟普通的 \texttt{"p"} 沒有兩樣。真正讓 Python 認出「這是 tuple」的關鍵，其實是那
個逗號，\texttt{("p",)} 裡的逗號才是重點，括號只是順便加上去、讓人看得更清楚而已。

\subsection*{簡單認識一個進階變體：namedtuple}

Python 標準函式庫裡有一個叫 \texttt{namedtuple} 的工具，可以幫 tuple 裡的每一個位置
額外取一個名字，讓你除了用 \texttt{coordinate[0]} 這種數字索引，也能用
\texttt{coordinate.latitude} 這種點加名字的方式取值，讀起來更清楚，也比較不容易搞錯
順序。

這個工具需要透過之後幾章才會教到的「把現成工具叫進來使用」的語法才能用，這裡先
不深入，只需要知道有這麼一個東西存在：它本質上還是一個 tuple（一樣不能修改內容），只
是多了一層「幫每個位置取名字」的方便包裝，等你之後真的碰到需要「一組固定欄位、但又想
用名字讀取」的情境，可以回頭來查它怎麼用。

\begin{drills}[label=drills:tuple-basics]
\item 這段程式碼會印出什麼？
\begin{verbatim}
value = ("hello")
print(type(value))
\end{verbatim}
\item 這段程式碼執行後會發生什麼事？寫出你預期看到的錯誤類型（不用寫出完整訊息）。
\begin{verbatim}
sizes = ("S", "M", "L")
sizes[1] = "XL"
\end{verbatim}
\item 這段程式碼執行後會發生什麼事？寫出你預期看到的錯誤類型，並說說看它跟上一題的錯誤
類型為什麼不一樣。
\begin{verbatim}
sizes = ("S", "M", "L")
sizes.append("XL")
\end{verbatim}
\item 這段程式碼會印出什麼？
\begin{verbatim}
weekday = ("Mon", "Tue")
weekend = ("Sat", "Sun")
print(weekday + weekend)
print(weekday * 2)
\end{verbatim}
\item 這段程式碼的執行結果會是 \texttt{True} 還是 \texttt{False}？
\begin{verbatim}
print((2, 9) < (2, 10))
\end{verbatim}
\end{drills}

\section{把 Tuple 跟賦值結合起來：Unpacking}
\label{sec:tuple-unpacking}

知道 tuple 是什麼之後，這一節要介紹這一章真正的重頭戲：把 tuple 跟賦值語法結合起來的
技巧，叫做 unpacking（解包）。概念很直接：如果右邊是一個裝著好幾個值的容器，左邊可以
一次列出好幾個變數名字，Python 會照順序把每個值分別存進對應的變數，一行程式碼就搞定，
不用再寫「先存成一個容器、再用索引一個一個挑出來」的繁瑣步驟。
\index{unpacking}

\begin{code}{python}{多重賦值的基本語法}
name1, name2, name3 = a_container_with_matching_length
\end{code}

\begin{example}{python}{(parse-date.py) 拆解日期字串}
date_string = "2026-09-01"
year, month, day = date_string.split("-")
print(f"Year: {year}, Month: {month}, Day: {day}")
\end{example}

\iosample{}{Year: 2026, Month: 09, Day: 01 \\}

這裡 \texttt{date\_string.split("-")} 會回傳一個裝著三個字串的 list（\texttt{["2026",
"09", "01"]}），unpacking 對 list 也一樣有效，不是只有 tuple 才能用，重點在「右邊的東
西裝了幾個值」，跟它是 tuple 還是 list 無關。跟原本要寫
\texttt{year = date\_string.split("-")[0]} 這種一個個挑出來的寫法比起來，unpacking 一
次把三個有意義的名字都準備好了，可讀性明顯好上不少。

\subsection*{常見錯誤：數量對不上}

unpacking 有一個硬性規定：左邊變數的個數，一定要跟右邊容器裡的元素個數完全一樣，多一
個少一個都不行。

\begin{example}{python}{(unpack-mismatch.py) 賦值數量對不上}
point = (10, 20, 30)
x, y = point
\end{example}

\begin{verbatim}
Traceback (most recent call last):
  File "unpack-mismatch.py", line 2, in <module>
    x, y = point
    ^^^^
ValueError: too many values to unpack (expected 2)
\end{verbatim}

\texttt{too many values to unpack} 是說「右邊的值太多了，左邊的變數裝不下」。反過來，
如果左邊變數比右邊值還多，會得到相反方向的抱怨：

\begin{verbatim}
Traceback (most recent call last):
  File "unpack-mismatch2.py", line 2, in <module>
    x, y, z = point
    ^^^^^^^
ValueError: not enough values to unpack (expected 3, got 2)
\end{verbatim}

兩則訊息都直接把「期望的個數」跟「實際拿到的個數」告訴你，通常一看就知道該加還是該減。

\subsection*{用底線忽略用不到的值}

有時候右邊容器裡的某個值你根本用不到，只是因為數量要對齊，不得不放一個變數名字接住它。
這種情況 Python 社群習慣用底線 \texttt{\_} 當變數名，代表「這個值我知道存在，但故意不
打算用它」。

\begin{example}{python}{(ignore-with-underscore.py) 用底線接住不需要的值}
student = ("Taylor", 16, "Denver")
name, _, city = student
print(f"{name} lives in {city}")
\end{example}

\iosample{}{Taylor lives in Denver \\}

\texttt{\_} 本身其實沒有任何特殊語法地位，它就是一個普通的變數名字，你甚至可以正常
印出 \texttt{print(\_)} 拿到那個被跳過的值。它會被拿來當「這裡故意不用」的訊號，純粹
是 Python 寫作者之間約定俗成的習慣，讓讀你程式碼的人一眼就看出「這個位置的值不重要，
不用花力氣去追它後面有沒有被用到」。

\subsection*{用星號一次接住「剩下的」值}

如果你只在乎前面幾個值，剩下不管有幾個都想整包收進一個 list 裡呢？在某個變數名前面加
一個星號 \texttt{*}，就能把「剩下的所有值」打包成一個 list，交給那個變數。

\begin{code}{python}{星號 unpacking 的三種常見寫法}
first, *rest = some_container      # 單獨接住第一個元素，其餘打包成 list
*init, last = some_container       # 單獨接住最後一個元素，前面打包成 list
first, *middle, last = some_container  # 接住第一個和最後一個，中間打包成 list
\end{code}

\begin{example}{python}{(star-unpacking.py) 用星號拆解一份成績清單}
scores = [95, 88, 76, 82, 91]

top, *others = scores
print(top)
print(others)

*rest, last = scores
print(rest)
print(last)

first, *middle, last = scores
print(first)
print(middle)
print(last)
\end{example}

\iosample{}{95 \\{} [88, 76, 82, 91] \\{} [95, 88, 76, 82] \\ 91 \\ 95 \\{} [88, 76, 82] \\ 91 \\}

留意每一次印出來的 \texttt{others}、\texttt{rest}、\texttt{middle} 都是一個 list，不
是 tuple。星號變數負責收留下來的所有值，數量不管是零個、一個還是很多個都沒問題，Python
會自動打包成一個 list 給你，不需要事先知道剩下確切有幾個值，這正是星號 unpacking 跟前
面「數量一定要對齊」規則最大的不同：一個星號變數可以吃下任意數量的值，讓左右兩邊不再需
要完全一樣長。

一個賦值裡最多只能出現一個星號變數，Python 沒辦法判斷兩個「剩下的」到底
該怎麼切分：

\begin{verbatim}
File "two-stars.py", line 1
    a, *b, *c = [1, 2, 3, 4]
    ^^^^^^^^^
SyntaxError: multiple starred expressions in assignment
\end{verbatim}

星號變數雖然可以吃下任意數量的值，但左邊固定位置的變數還是至少要有值可以對應，如果右邊
的元素數量連固定位置都湊不滿，一樣會噴出 \texttt{ValueError}：

\begin{verbatim}
Traceback (most recent call last):
  File "not-enough-for-star.py", line 1, in <module>
    a, b, *rest = [1]
    ^^^^^^^^^^^
ValueError: not enough values to unpack (expected at least 2, got 1)
\end{verbatim}

留意這裡的措辭是 \texttt{expected at least 2}，跟前面沒有星號時的
\texttt{expected 3, got 2} 不太一樣：因為星號變數願意接受「零個或更多」，Python 只要
求固定位置的變數（這裡是 \texttt{a}、\texttt{b}）一定要湊得滿，星號那個位置本身沒有下
限。

星號 unpacking 在處理「第一列是標題、後面每一列才是真正資料」這種常見的資料格式時特別
好用，之後第 \ref{ch:csv} 章處理 CSV 檔案時會常常遇到這種形狀的資料。

\begin{example}{python}{(skip-header-row.py) 用星號拆出標題與其餘資料列}
raw_data = [
    "Name,Score",
    "Taylor,95",
    "Jordan,88",
]

header, *records = raw_data
print(header)
print(records)
\end{example}

\iosample{}{Name,Score \\{} ['Taylor,95', 'Jordan,88'] \\}

\texttt{header} 接住第一列標題，\texttt{*records} 把剩下不管幾列，全部打包成一個
list，不需要事先知道這份資料到底有幾列，這正是星號 unpacking 比「數清楚長度、再用切片
手動拆」方便的地方。

\subsection*{用 unpacking 交換兩個變數}

unpacking 有一個很受歡迎的小應用：交換兩個變數的值。

\begin{example}{python}{(swap.py) 交換兩個變數的值}
a = "mango"
b = "lemon"
print(f"Before swap: a = {a}, b = {b}")

a, b = b, a

print(f"After swap: a = {a}, b = {b}")
\end{example}

\iosample{}{Before swap: a = mango, b = lemon \\ After swap: a = lemon, b = mango \\}

\stylepoint{很多語言要交換兩個變數的值，得先借一個第三個暫存變數：把 \texttt{a} 存到暫
存變數，再把 \texttt{b} 存進 \texttt{a}，最後把暫存變數存進 \texttt{b}。Python 直接一
行 \texttt{a, b = b, a} 就做完了。右邊的 \texttt{b, a} 先被組成一個 tuple，記住兩個舊
值，再一次 unpack 回左邊的 \texttt{a} 和 \texttt{b}，完全不需要暫存變數，這是很多人覺
得「這樣寫比較 Python」的經典例子。}

這個技巧不只限於兩個變數，三個（或更多）變數要「輪流交換」時，同樣的寫法一樣管用。

\begin{example}{python}{(rotate-three.py) 三個變數輪流交換}
a = 1
b = 2
c = 3
print(f"Before swap: a = {a}, b = {b}, c = {c}")

a, b, c = c, a, b

print(f"After swap: a = {a}, b = {b}, c = {c}")
\end{example}

\iosample{}{Before swap: a = 1, b = 2, c = 3 \\ After swap: a = 3, b = 1, c = 2 \\}

原理跟兩個變數的 swap 完全一樣：右邊的 \texttt{c, a, b} 先被組成一個 tuple，記住三個舊
值，再一次 unpack 回左邊的 \texttt{a}、\texttt{b}、\texttt{c}，每個變數拿到的都是「上
一個位置原本的值」，一次到位，不需要暫存變數，也不需要拆成三行分開賦值。

\subsection*{\texttt{for} 迴圈中的 unpacking}

第 \ref{ch:dict-set} 章你學過用 \texttt{.items()} 走訪字典的 key 跟 value，當時的寫法
\texttt{for key, value in some\_dict.items():} 其實就是 unpacking，\texttt{.items()}
每一輪給你的都是一個 \texttt{(key, value)} 的 tuple，\texttt{for} 迴圈把它直接 unpack
成兩個變數，原來那個語法背後就是這一章教的東西。

\begin{example}{python}{(dict-items-unpacking.py) 走訪字典的 key 與 value}
prices = {"apple": 35, "banana": 20, "orange": 45}

for name, price in prices.items():
    print(f"{name} costs {price} dollars each")
\end{example}

\iosample{}{apple costs 35 dollars each \\ banana costs 20 dollars each \\ orange costs 45 dollars each \\}

\subsection*{巢狀 unpacking：tuple 裡面還裝著 tuple}

如果容器裡的元素本身又是一個 tuple，unpacking 的寫法也可以跟著「巢狀」，只要讓左邊的
形狀跟右邊對得上，Python 就能一路拆到底。

\begin{example}{python}{(nested-unpacking.py) 走訪一份「城市配座標」的清單}
points = [
    ("Boston", (25.03, 121.56)),
    ("Chicago", (24.15, 120.67)),
]

for city, (lat, lon) in points:
    print(f"{city}: latitude {lat}, longitude {lon}")
\end{example}

\iosample{}{Boston: latitude 25.03, longitude 121.56 \\ Chicago: latitude 24.15, longitude 120.67 \\}

\texttt{points} 裡的每一個元素，都是一個 \texttt{(城市名稱, (緯度, 經度))} 這樣兩層的
tuple。\texttt{for} 迴圈的 \texttt{city, (lat, lon)} 這個寫法，形狀跟右邊完全對應：外
層的兩個位置分別對到 \texttt{city} 跟一整個座標 tuple，而座標 tuple 本身又立刻被拆成
\texttt{lat}、\texttt{lon} 兩個變數。左右兩邊的「形狀」要吻合，Python 才知道該怎麼拆，
這跟第 1 節提到的「數量要對齊」是同一個規則，只是這裡多了一層。

\subsection*{函式同時回傳多個值}

第 \ref{ch:functions} 章你學過函式用 \texttt{return} 回傳一個結果。如果想讓函式同時給出好幾個結果呢？
做法是讓 \texttt{return} 後面接好幾個用逗號分隔的值，這其實是在回傳一個 tuple，函式本
質上還是只能回傳一個東西，只是這個東西剛好是一個 tuple。呼叫端再用 unpacking 把它拆開，
分別存進有意義名字的變數。

\begin{example}{python}{(min-max.py) 一次算出最大值和最小值}
def min_max(numbers):
    return min(numbers), max(numbers)


data = [8, 3, 15, 1, 9]
lo, hi = min_max(data)
print(f"Min: {lo}, Max: {hi}")
\end{example}

\iosample{}{Min: 1, Max: 15 \\}

如果呼叫端忘記 unpack，只用一個變數接住呼叫結果，會發生什麼事？並不會出錯，只是你拿到
的會是整個 tuple 本身，沒有被拆開。

\begin{example}{python}{(forget-to-unpack.py) 忘記 unpack 會拿到整個 tuple}
def min_max(numbers):
    return min(numbers), max(numbers)


data = [8, 3, 15, 1, 9]
result = min_max(data)
print(result)
print(type(result))
\end{example}

\iosample{}{(1, 15) \\{} <class 'tuple'> \\}

這正好提醒你，函式回傳多個值這件事，從函式的角度看，它從頭到尾都只回傳了一個東西
（一個 tuple），「同時拿到兩個變數」是呼叫端 unpack 之後才有的結果，不是函式本身的
魔法。忘記 unpack 不會報錯，但你接下來想用 \texttt{lo}、\texttt{hi} 這兩個名字的地方，
就會因為它們根本不存在而噴出 \texttt{NameError}。

其實不用等到你自己寫函式，Python 內建函式裡就有現成的例子。\texttt{divmod(a, b)} 會
一次算出 \texttt{a} 除以 \texttt{b} 的商跟餘數，回傳的正是一個裝著兩個值的 tuple。

\begin{example}{python}{(divmod-example.py) 用內建函式 divmod() 一次拿到商與餘數}
quotient, remainder = divmod(17, 5)
print(f"Quotient is {quotient}, remainder is {remainder}")
\end{example}

\iosample{}{Quotient is 3, remainder is 2 \\}

這正好說明「函式回傳多個值」不是什麼特殊魔法，只要函式回傳的是一個 tuple，呼叫端就能
用這一節學到的 unpacking 語法把它拆開，不管這個函式是你自己寫的，還是 Python 內建的。

\begin{drills}[label=drills:tuple-unpacking]
\item 這段程式碼會印出什麼？
\begin{verbatim}
first = "cat"
second = "dog"
first, second = second, first
print(first, second)
\end{verbatim}
\item 這段程式碼會印出什麼？
\begin{verbatim}
inventory = {"notebook": 3, "pen": 12}
for item, count in inventory.items():
    print(f"{item}: {count}")
\end{verbatim}
\item 這段程式碼會噴出 \texttt{ValueError}。要怎麼修改（不改變 \texttt{record} 那一
行），才能讓它正常執行？
\begin{verbatim}
record = ("Alice", 15, "Boston")
name, age = record
\end{verbatim}
\item 一個函式定義如下，呼叫端要怎麼寫，才能同時拿到 \texttt{total} 跟
\texttt{average} 兩個變數？
\begin{verbatim}
def total_and_average(numbers):
    return sum(numbers), sum(numbers) / len(numbers)
\end{verbatim}
\item 這段程式碼會印出什麼？
\begin{verbatim}
values = [10, 20, 30, 40]
first, *rest = values
print(rest)
\end{verbatim}
\item 這段程式碼會印出什麼？
\begin{verbatim}
pair = ("Boston Station", (25.05, 121.52))
name, (lat, lon) = pair
print(name, lat)
\end{verbatim}
\item 這段程式碼會印出什麼？
\begin{verbatim}
a, b, c = 1, 2, 3
a, b, c = b, c, a
print(a, b, c)
\end{verbatim}
\end{drills}

\section{不可變 vs 可變：名字牌模型的收尾對比}
\label{sec:tuple-vs-list}

還記得第 1 節撞牆的時候留下的那個問題嗎？如果「等號是幫物件貼名字牌，不是複製」這件事
是真的，那為什麼 tuple 看起來「整個換掉」的時候，感覺就像被複製了一份新的一樣？這一節
就要把這件事講清楚。

先看一段程式碼，雖然 tuple 不能被原地修改，但你還是可以寫出「看起來像修改」的程式碼，
例如把某個位置的值換掉、拼接一個新的 tuple 上去：

\begin{example}{python}{(replace-in-tuple.py) 看起來像修改，其實是造新物件}
scores = ("A", "B", "C")
print(scores, id(scores))

scores = ("Z",) + scores[1:]
print(scores, id(scores))
\end{example}

\iosample{}{('A', 'B', 'C') 4349906432 \\ ('Z', 'B', 'C') 4349912128 \\}

（你自己實際跑出來的那兩個數字，幾乎一定跟這裡不一樣，\texttt{id()} 回傳的是物件在記
憶體裡的位址，每次執行、每台機器都可能不同。重點不是數字本身，是\textbf{前後兩個數字不
一樣}。）

\texttt{id()} 是 Python 內建的函式，能告訴你一個物件在記憶體裡的「門牌號碼」，兩個物
件的 \texttt{id()} 一樣，就代表它們是同一個物件。\index{id()}上面這段程式碼裡，\texttt{scores}
「看起來」被改成了 \texttt{('Z', 'B', 'C')}，但 \texttt{id()} 告訴你，第二行印出來的
\texttt{scores} 已經是一個全新的物件了，跟第一行那個完全不是同一個東西。

這正好回答了開頭的問題：\textbf{不是名字牌模型在 tuple 這裡失效了，是 tuple 不可變，所
以每一次「看起來像修改」的操作，都逼著 Python 造一個全新的物件，再把 \texttt{scores}
這張名字牌重新貼上去}。第二行的 \texttt{scores = ("Z",) + scores[1:]} 其實分兩步：先用
\texttt{("Z",) + scores[1:]} 造出一個全新的 tuple，再把 \texttt{scores} 這張名字牌從舊
物件撕下來、貼到這個新物件上，原本那個 \texttt{('A', 'B', 'C')} 並沒有被改變，它只是
不再被任何名字牌指著而已。

拿第 \ref{ch:lists} 章 list 的別名行為回頭對照，會看得更清楚。如果兩張名字牌指著同一個 list，透過其
中一個改動內容，另一個看到的也是改過的結果，因為 list 可以被原地修改，兩張名字牌指的
從頭到尾都是同一個物件。

但如果換成 tuple，同樣「兩張名字牌指著同一個物件」的情況，你不
可能用某張名字牌去「原地修改」它，因為 tuple 根本沒有能原地修改自己的操作，任何看起來
像修改的動作，都只能是造一個新物件、把其中一張名字牌重新貼上去，另一張名字牌仍然貼在原
本那個沒被動過的舊物件上。

\keypoint{對 int、字串、tuple 這類不可變型別，「名字牌模型」跟你原本從數學課或日常生活
帶來的「盒子」直覺，在這裡剛好會給出一樣的結果，因為東西不能被原地修改，每次「看起來
像修改」都只能整個換掉，感覺就很像「盒子裡的內容被換掉了」。但這只是巧合，不是名字牌模
型的例外：底層規則從頭到尾都沒有變，變的只是「這個型別到底允不允許被原地修改」。}

\subsection*{一個更細微的對照：\texttt{+=} 到底有沒有「原地修改」}

第 \ref{ch:loops} 章你學過 \texttt{+=} 是「把右邊的值加到變數原本的值上，再把結果重新貼回同一個名
字牌」的簡寫。這句話聽起來好像 \texttt{+=} 一定會造出新物件，但事情沒有這麼絕對：對
list 使用 \texttt{+=}，Python 其實會直接在原本那個物件上把新內容接上去，物件本身沒有
換過；對 tuple 使用 \texttt{+=}，因為 tuple 沒有任何原地修改自己的操作，只能乖乖造一
個新物件。用 \texttt{id()} 可以親眼看到這個差別。

\begin{example}{python}{(plus-equals-tuple-vs-list.py) 用 id() 比較 += 對 list 和 tuple 的差別}
numbers_list = [1, 2, 3]
before_list_id = id(numbers_list)
numbers_list += [4]
print(numbers_list, id(numbers_list) == before_list_id)

numbers_tuple = (1, 2, 3)
before_tuple_id = id(numbers_tuple)
numbers_tuple += (4,)
print(numbers_tuple, id(numbers_tuple) == before_tuple_id)
\end{example}

\iosample{}{[1, 2, 3, 4] True \\ (1, 2, 3, 4) False \\}

第一行的 \texttt{True} 告訴你，\texttt{numbers\_list} 用 \texttt{+=} 之後，還是原本
那個物件，只是內容被原地擴充了；第二行的 \texttt{False} 告訴你，\texttt{numbers\_tuple}
用 \texttt{+=} 之後，已經是一個全新的物件了。同一個運算子、幾乎一模一樣的寫法，因為
兩邊型別的可變性不同，骨子裡做的事情完全不一樣：這正是這一章從頭到尾一直在強調的重點，
「型別允不允許原地修改」才是決定 Python 背後到底做了什麼事的關鍵，不是看程式碼表面長
得像不像。

\subsection*{tuple 不可變，指的是「這一層」不能換，不是裡面的東西永遠凍結}

這裡有一個很容易被誤會的細節：tuple 保證的是「這個 tuple 本身裝著哪些物件」不能被改
變，但如果裡面裝的某個元素本身是可變的（例如一個 list），那個元素自己還是可以被原地修
改，跟外層的 tuple 完全是兩回事。

\begin{example}{python}{(mutable-inside-tuple.py) tuple 裡面的 list 還是可以被修改}
data = (1, 2, [3, 4])
print(data)

data[2].append(5)
print(data)
\end{example}

\iosample{}{(1, 2, [3, 4]) \\ (1, 2, [3, 4, 5]) \\}

\texttt{data[2] = "something"} 這種「把 tuple 裡第三個位置整個換掉」的寫法一樣會踩到
第 1 節那道牆，觸發 \texttt{TypeError}。但 \texttt{data[2].append(5)} 完全是另一回事：
它沒有動 \texttt{data} 這個 tuple 裡「第三個位置放的是哪個物件」這件事，第三個位置從
頭到尾都指著同一個 list，只是那個 list 自己被原地修改了內容。tuple 只管住自己這一層，
管不到（也沒打算管）裡面的物件要不要允許被修改。

這個細節也解釋了為什麼 hashable（可雜湊）的判斷要看得更仔細一點：如果一個 tuple 裡裝
了像 list 這種不可雜湊的東西，這個 tuple 整體也會變成不可雜湊，一樣不能拿來當 dict 的
key。

\begin{verbatim}
Traceback (most recent call last):
  File "mutable-inside-tuple.py", line 3, in <module>
    d = {data: "some value"}
        ^^^^^^^^^^^^^^^^^^^^
TypeError: unhashable type: 'list'
\end{verbatim}

錯誤訊息抱怨的是 \texttt{'list'} 而不是 \texttt{'tuple'}，因為 Python 檢查一個 tuple
能不能被雜湊時，會連著檢查它裡面每一個元素能不能被雜湊，只要有一個不行，整個 tuple 就
不行。

\keypoint{tuple 保證的不可變是「淺層」的：只保證它自己這一層的內容（裝著哪些物件）不會
被換掉，不保證每個裝進去的物件本身也不能被修改。想要一份資料整個都不能被動到，裡面每
一層也都要用不可變型別組成，例如全部用 tuple、字串、數字組成，不能有任何一層是 list
或 dict。}

\subsection*{不可變的實際好處：tuple 可以當 dict 的 key}

「內容不會變」這件事，除了避免你意外改壞資料，還有一個很實際的好處：tuple 可以拿來當
dict 的 key，而 list 不行。

\begin{example}{python}{(tuple-as-dict-key.py) 用座標當字典的 key}
locations = {
    (25.03, 121.56): "Boston Tower",
    (22.63, 120.30): "Seattle Port",
}
print(locations[(25.03, 121.56)])
\end{example}

\iosample{}{Boston Tower \\}

如果改用 list 當 key，Python 會直接拒絕：

\begin{verbatim}
Traceback (most recent call last):
  File "tuple-as-dict-key.py", line 2, in <module>
    locations = {bad_key: "Boston Tower"}
                ^^^^^^^^^^^^^^^^^^^^^
TypeError: unhashable type: 'list'
\end{verbatim}

\texttt{unhashable} 的意思是「沒辦法被雜湊」，dict 的 key 內部依賴一個叫「雜湊值」的
機制才能快速查找，而這個機制要求 key 的內容永遠不會變，這樣雜湊值才不會失效。list 的內
容隨時可能被改，Python 乾脆不允許它當 key；tuple 造出來就不能再改，正好滿足這個要求，
所以可以放心拿來當 key。\index{hashable}這也是前面提到的 \texttt{namedtuple} 一個常被拿來用的場景：既
要一組固定的欄位當 key，又想讓讀程式碼的人看得出每個位置代表什麼意思。

\begin{drills}[label=drills:tuple-vs-list]
\item 這段程式碼執行完，\texttt{a is b} 會是 \texttt{True} 還是 \texttt{False}？
\begin{verbatim}
a = [1, 2, 3]
b = a
b.append(4)
\end{verbatim}
\item 這段程式碼執行完，\texttt{t is t2} 會是 \texttt{True} 還是 \texttt{False}？
\begin{verbatim}
t = (1, 2, 3)
t2 = t
t2 = t2 + (4,)
\end{verbatim}
\item 承上題，如果把最後一行改成 \texttt{t2 = t}（不做任何加法），\texttt{t is t2}
還會是同樣的答案嗎？為什麼？
\item 這段程式碼執行完，\texttt{original\_id == id(nums)} 會是 \texttt{True} 還是
\texttt{False}？
\begin{verbatim}
nums = [1, 2, 3]
original_id = id(nums)
nums += [4, 5]
\end{verbatim}
\item 這段程式碼執行完，\texttt{t} 的內容會是什麼？tuple 本身有沒有被「修改」？
\begin{verbatim}
t = (1, [2, 3])
t[1].append(4)
\end{verbatim}
\end{drills}

\section*{練習解答}
\addcontentsline{toc}{section}{練習解答}

\paragraph{第~\ref{sec:tuple-basics}~節}
\begin{enumerate}
\item 印出 \texttt{<class 'str'>}，單獨一個括號包住的值不會變成 tuple，少了逗號，
Python 只當它是一個字串。
\item \texttt{TypeError}，tuple 不支援用索引賦值來改變內容
（\texttt{'tuple' object does not support item assignment}）。
\item \texttt{AttributeError}，因為 tuple 從一開始就沒有 \texttt{.append()} 這個方
法，跟上一題「有這個操作但不准做」的 \texttt{TypeError} 不一樣，這裡是「這個操作根本不
存在」。
\item 印出：
\begin{verbatim}
('Mon', 'Tue', 'Sat', 'Sun')
('Mon', 'Tue', 'Mon', 'Tue')
\end{verbatim}
\texttt{+} 把兩個 tuple 接成一個新 tuple，\texttt{*} 把整份內容重複兩次。
\item \texttt{True}，兩個 tuple 第一個元素都是 \texttt{2}，分不出高下，往下比第二個
元素，\texttt{9 < 10} 成立，比較到這裡就有答案了。
\end{enumerate}

\paragraph{第~\ref{sec:tuple-unpacking}~節}
\begin{enumerate}
\item 印出 \texttt{dog cat}，\texttt{first, second = second, first} 先把 \texttt{dog}、
\texttt{cat} 組成一個暫時的 tuple，再 unpack 回去，達成交換。
\item 印出：
\begin{verbatim}
notebook: 3
pen: 12
\end{verbatim}
\item \texttt{record} 有三個元素，但左邊只列了兩個變數，數量對不上。修法是讓左邊也列
出三個變數，例如 \texttt{name, age, city = record}（如果不需要用到 \texttt{city}，也
可以用一個底線當佔位名字：\texttt{name, age, \_ = record}）。
\item \texttt{total, average = total\_and\_average(某個數字 list)}，函式回傳的是一個
裝著兩個值的 tuple，呼叫端用 unpacking 把它拆成兩個變數接住。
\item 印出 \texttt{[20, 30, 40]}，星號變數 \texttt{*rest} 會把除了第一個以外剩下的所
有值打包成一個 list。
\item 印出 \texttt{Boston Station 25.05}，\texttt{name} 對到外層第一個位置，\texttt{(lat, lon)}
這個形狀對到外層第二個位置（本身又是一個座標 tuple），再拆成 \texttt{lat}、
\texttt{lon} 兩個變數，這裡只印出 \texttt{name} 和 \texttt{lat}。
\item 印出 \texttt{2 3 1}，右邊的 \texttt{b, c, a} 先依照舊值組成一個 tuple
\texttt{(2, 3, 1)}，再一次 unpack 回左邊的 \texttt{a}、\texttt{b}、\texttt{c}，等於
把三個值往前輪轉一格。
\end{enumerate}

\paragraph{第~\ref{sec:tuple-vs-list}~節}
\begin{enumerate}
\item \texttt{True}，\texttt{b = a} 只是幫同一個 list 多貼一張名字牌，\texttt{.append()}
是原地修改，兩張名字牌指的還是同一個物件。
\item \texttt{False}，\texttt{t2 = t2 + (4,)} 造出一個全新的 tuple，再把 \texttt{t2}
這張名字牌重新貼到新物件上，\texttt{t} 仍然貼在原本沒被動過的舊物件上。
\item 不會，答案會變成 \texttt{True}，這次沒有做任何「看起來像修改」的操作，
\texttt{t2 = t} 單純是幫同一個物件多貼一張名字牌 \texttt{t2}，跟 list 的別名行為（第一
題）道理完全一樣，只是這裡剛好用在 tuple 上而已。
\item \texttt{True}，對 list 使用 \texttt{+=} 是原地修改，\texttt{nums} 這張名字牌從
頭到尾都貼在同一個物件上，跟前一節 \texttt{plus-equals-tuple-vs-list.py} 範例裡 list
那一半的結果一致。
\item \texttt{t} 會變成 \texttt{(1, [2, 3, 4])}；tuple 本身沒有被修改，它從頭到尾都是
同一個物件、一樣裝著兩個位置，只是第二個位置指著的那個 list 被原地修改了內容。
\end{enumerate}
